\documentclass[10pt,a4paper]{amsart}
\usepackage[margin=19mm]{geometry}
\usepackage{amsmath,amssymb,mathtools}
\usepackage[scr=rsfs,cal=euler]{mathalfa}
\usepackage{xfrac}
\usepackage[unicode,hidelinks,bookmarks=false]{hyperref}
\newcommand{\ie}{i.e.\ }
\newcommand{\cf}{cf.\ }
\newcommand{\resp}{respectively\ }
\newcommand{\Subsection}{\subsection}
\providecommand{\nobreakdash}{\nobreak\hbox{-}}
\setlength{\emergencystretch}{2em}
\multlinegap0pt
\usepackage{bm}
\usepackage{bbm}
\usepackage{dsfont}
\usepackage{xspace}
\usepackage{cancel}
\usepackage{mathtools}
\usepackage{amsthm}
\makeatletter
\@ifundefined{theorem}{\newtheorem{theorem}{Theorem}}{}
\makeatother
\title[A Unified Lyapunov-IQC Framework for Uniform Stability of Sm]{A Unified Lyapunov-IQC Framework for Uniform Stability of Smooth Quadratic First-Order Accelerated Optimizers}
\author{Don Li, Dacian Daescu}
\date{Source: arXiv:2605.08488v2 (CC BY 4.0)}
\begin{document}
\maketitle

\noindent\emph{Source and licence.} A Unified Lyapunov-IQC Framework for Uniform Stability of Smooth Quadratic First-Order Accelerated Optimizers. Version arXiv:2605.08488v2 (CC BY 4.0), distributed under the Creative Commons Attribution 4.0 International licence, as recorded in the arXiv licence metadata for this exact version and shown on the abstract page https://arxiv.org/abs/2605.08488v2. Full rights notice: licenses/arxiv-2605.08488-CC-BY-4.0.txt.\par\smallskip

\noindent\textit{Editorial scope.} Counted unit: the Theorem whose source label is \texttt{None} in the article (source0.tex lines 1204--1210, source label \texttt{None}), delivered together with the article's own complete original proof of it (source0.tex lines 1212--1390, one \texttt{proof} environment of 179 lines). No mathematics is written, repaired or re-derived: statement and proof are the article's own text line for line. Required context retained uncounted: none. Every definition, display and label of those lines is retained unchanged. Omitted from this package and not redistributed: 1--1203, 1211--1211, 1391--1393. Nothing inside a retained excerpt is omitted or altered: every retained body is delivered exactly as the article's lines are written, so no omission receipt of any kind accompanies this package. Citations: the retained excerpts carry no citation command at all, so no bibliography entry of the article is transcribed and no reference is redistributed. Reference adaptation: none was needed and none was made; every reference inside the delivered unit resolves inside it, so no unresolved reference or citation is produced. Printed slips kept literal: no printed word, symbol, hypothesis or number of the counted statement or of its proof was altered. Layout and class: the article's own class is \texttt{article} with packages including jmlr2e, algpseudocode, algorithm, float, amsmath, amssymb, optidef; the wrapper is the lane's pinned \texttt{amsart} editorial wrapper at 10pt on A4, which restates the theorem-like environments the retained text needs and adds the article's own macro definitions that the retained text needs (none), with extra wrapper packages (bm, bbm, dsfont, xspace, cancel, mathtools). The delivered numbering is the unit's own. Assets: no retained span contains a figure environment, a graphics-inclusion command, a PDF-page inclusion, a table or a TikZ picture; no image file is copied into this package, the package declares no asset dependency and no image file is redistributed with it. Licence: version arXiv:2605.08488v2 of this article is distributed under the Creative Commons Attribution 4.0 International licence, as recorded in the arXiv licence metadata for this exact version and shown on the abstract page https://arxiv.org/abs/2605.08488v2. A full rights notice is delivered at licenses/arxiv-2605.08488-CC-BY-4.0.txt. Characters: every retained excerpt is pure ASCII, so the delivered engine, XeLaTeX with its default Latin Modern text font, reports no missing character.
\begin{theorem}
Assume that the loss $f_i(w)$ is $\gamma$-strongly convex and $\beta$-smooth for all $z$. Run SGD with constant $\eta \leq \frac{\gamma}{\beta^2}$ for $T$ iterations. Then SGD is uniformly stable with

\begin{center}
$\epsilon \leq \frac{2G^2}{\gamma n}$.
\end{center}
\end{theorem}

\begin{proof}
Let $S$ and $S'$ be two samples of size $n$ that differ only at a single point. We want to show that SGD is $\epsilon$-uniformly stable, i.e., by definition,

\begin{center}
$\text{sup}_z \mathbb{E}[\ell(A(S), z) - \ell(A(S'), z)] \leq \epsilon$
\end{center}

\noindent
for arbitrary point $z$. Let SGD run over coupled indices over $S$ and $S'$, where $S = (z_1, z_2, ... , z_n)$ and $S' = (z_1', z_2', ... , z_n')$. We denote the index of the differing data point as $j \in \{1,2, ... , n\}$, i.e., $z_j \neq z_j'$ (so $z_i = z_i'$ for all $i \neq j$). \\
\\
Recall the update rule for SGD for some arbitrary parameter $x_t$, i.e., 

\begin{center}
$x_{t+1} = x_t - \eta \nabla f_i(x_t)$,
\end{center}

\noindent
where $\eta$ is the learning rate/step size and $\nabla f_i(x_t)$ is the gradient at $x_t$. Run SGD over both $S$ and $S'$. In doing so, we obtain two separate update rules:

\begin{equation}
w_{t+1} = w_t - \eta \nabla f_i(w_t),
\end{equation}

\begin{equation}
w'_{t+1} = w'_t - \eta \nabla f_i(w'_t).
\end{equation}

\noindent
We now define an iterate divergence term,

\begin{equation}
\delta_t = \mathbb{E}[||w_t - w'_t||].
\end{equation}

\noindent
By definition, SGD is stable iff $\delta_t = O(\frac{1}{n})$. Note that there are two possible cases, (1) where the differing point is \textit{not} selected, i.e., where $i_t \neq j$, and (2) where the differing point \textit{is} selected, i.e., where $i_t = j$. Thus, we use proof by cases. \\
\\
\fbox{\{$i_t \neq j$\}} Let $\{z_1, z_2, ... , z_n\}$ be uniformly distributed, so $\mathbb{P}[i_t \neq j] = 1 - \frac{1}{n}$. Then we obtain Eqns. 3 \& 4 as the SGD update rules with the loss function and gradient values. Thus, $\delta_{t+1} = ||w_{t+1} - w'_{t+1}|| = ||(w_t - w'_t) - \eta(\nabla \ell(w_t, z_{i_t}) - \nabla \ell(w'_t, z'_{i_t}))||$. This implies $||w_{t+1} - w'_{t+1}||^2 = ||(w_t - w'_t) - \eta(\nabla \ell(w_t, z_{i_t}) - \nabla \ell(w'_t, z'_{i_t}))||^2 = || w_{t} - w'_{t}||^2 - 2\eta \langle w_t - w'_t, \nabla \ell(w_t, z_{i_t}) - \nabla \ell(w'_t, z'_{i_t} \rangle + \eta^2|| \nabla \ell(w_t, z_{i_t}) - \nabla \ell(w'_t, z'_{i_t})||^2$, i.e., 

\begin{equation}
||w_{t+1} - w'_{t+1}||^2 = || w_{t} - w'_{t}||^2 -2\eta \langle w_t - w'_t, \nabla \ell(w_t, z_{i_t}) - \nabla \ell(w'_t, z'_{i_t} \rangle + \eta^2|| \nabla \ell(w_t, z_{i_t}) - \nabla \ell(w'_t, z'_{i_t})||^2.
\end{equation}

\noindent
We bound the middle term, $-2\eta \langle w_t - w'_t, \nabla \ell(w_t, z_{i_t}) - \nabla \ell(w'_t, z'_{i_t} \rangle$, which we do via gradient monotonicity. Doing so yields

\begin{equation}
-2\eta \langle w_t - w'_t, \nabla \ell(w_t, z_{i_t}) - \nabla \ell(w'_t, z'_{i_t}) \rangle \leq -2\eta\gamma || w_t - w_t' ||^2.
\end{equation}

\noindent
We can now bound the last term, $\eta^2|| \nabla \ell(w_t, z_{i_t}) - \nabla \ell(w'_t, z'_{i_t})||^2$, which we do via $\beta$-smoothness. Directly applying the definition of $\beta$-smoothness yields

\begin{equation}
\eta^2|| \nabla \ell(w_t, z_{i_t}) - \nabla \ell(w'_t, z'_{i_t})||^2 \leq \eta^2 \beta^2|| w_t - w'_t ||^2.
\end{equation}

\noindent
These inequalities thus yield

\begin{equation}
||w_{t+1} - w'_{t+1}||^2 \leq ||w_{t} - w'_{t}||^2 - 2\eta\gamma|| w_t - w_t' ||^2 + \eta^2 \beta^2|| w_t - w'_t ||^2.
\end{equation}

\noindent
Factoring $||w_{t} - w'_{t}||^2$ on the right-hand side of the inequality yields

\begin{equation}
||w_{t+1} - w'_{t+1}||^2 \leq (1-2\eta\gamma+\eta^2\beta^2)||w_{t} - w'_{t}||^2.
\end{equation}

\noindent
Optimal $\eta$ occurs where $\frac{\delta}{\delta \eta}[1-2\eta\gamma+\eta^2\beta^2]=0$, which suggests $\eta \leq \frac{\gamma}{\beta^2}$. Then for $\eta \leq \frac{\gamma}{\beta^2}$ we have 

\begin{equation}
||w_{t+1} - w'_{t+1}||^2 \leq (1-\eta\gamma)||w_{t} - w'_{t}||^2.
\end{equation}

\noindent
Then taking the square root with a looser bound yields

\begin{equation}
||w_{t+1} - w'_{t+1}|| \leq (1-\eta\gamma)||w_{t} - w'_{t}||.
\end{equation}

\noindent
Thus, in the case where $\{i_t \neq j\}$, the iterate divergence decreases from one iteration to the next at a factor of $1 - \eta\gamma$. \\
\\
\fbox{$\{i_t = j\}$} We keep the same uniform distribution assumption for both $S$ and $S'$, where SGD is ran over a coupled set of indices over both $S$ and $S'$, so $\mathbb{P}[i_t = j] = \frac{1}{n}$. Recall $\delta_{t+1} = ||w_{t+1} - w'_{t+1}|| = ||(w_t - w'_t) - \eta(\nabla \ell(w_t, z_{i_t}) - \nabla \ell(w'_t, z'_{i_t}))||$. Applying the triangle inequality implies

\begin{equation}
||w_{t+1} - w'_{t+1}|| \leq ||(w_t - w'_t)|| +  \eta(\nabla \ell(w_t, z_{i_t}) - \nabla \ell(w'_t, z'_{i_t}))||.
\end{equation}

\noindent
We can then apply the triangle inequality again to the $(\nabla \ell(w_t, z_{i_t}) - \nabla \ell(w'_t, z'_{i_t}))||$, which yields

\begin{equation}
||w_{t+1} - w'_{t+1}|| \leq ||(w_t - w'_t)|| +  \eta||\nabla \ell(w_t, z_{i_t})|| + \eta ||\nabla \ell(w'_t, z'_{i_t}))||.  
\end{equation}

\noindent
By boundedness of the gradient, we can infer $||\nabla \ell(w_t, z_{i_t})|| \leq G$ and $||\nabla \ell(w'_t, z'_{i_t}))|| \leq G$, which implies

\begin{equation}
||w_{t+1} - w'_{t+1}|| \leq ||(w_t - w'_t)|| +  \eta G + \eta G = ||(w_t - w'_t)|| + 2\eta G.  
\end{equation}

\noindent
Let $\mathbb{E}[\delta_t] = ||(w_t - w'_t)||$. Taking expectation over (100) yields

\begin{equation}
\mathbb{E}[||w_{t+1} - w'_{t+1}|| | i_t = j] \leq \delta_t + 2 \eta G.
\end{equation}

\noindent
Moreover, we let $\delta_{t+1} = \mathbb{E}[||w_{t+1} - w'_{t+1}|| | i_t = j]$. Thus, for $i_t = j$, 

\begin{equation}
\delta_{t+1} \leq \delta_t + 2 \eta G.
\end{equation}

\noindent
Recall the law of total expectation for a discrete random variable $X$ with probability $P(A)$ for event $A$, where $A^c$ is the complement of $A$ with $P(X^c) = P(A^c)$, i.e., $\mathbb{E}[X] = P(A)\mathbb{E}[X|A] + P(A^c)\mathbb{E}[X|A^c]$. Here, $A$ is the event where $i_t \neq j$, so $A^c$ is the event where $i_t = j$. Thus, applying the law of total expectation yields

\begin{equation}
\delta_{t+1} = (1-\frac{1}{n})\mathbb{E}[||w_{t+1} - w'_{t+1}|| | i_t \neq j] + \frac{1}{n}\mathbb{E}[||w_{t+1} - w'_{t+1}|| | i_t = j].
\end{equation}

\noindent
We derived $\mathbb{E}[||w_{t+1} - w'_{t+1}|| | i_t \neq j] \leq (1-\eta\gamma)\delta_t$ and $\mathbb{E}[||w_{t+1} - w'_{t+1}|| | i_t = j] \leq \delta_t + 2\eta G$. Substituting this yields

\begin{equation}
\delta_{t+1} \leq (1-\frac{1}{n})(1-\eta\gamma)\delta_t + (\frac{1}{n})(\delta_t + 2\eta G).
\end{equation}

\noindent
Expansion of the right-hand side yields

\begin{equation}
\delta_{t+1} \leq (1 - \eta\gamma - \frac{1}{n} - \frac{\eta \gamma}{n})\delta_t + (\frac{1}{n})(\delta_t + 2\eta G),
\end{equation}

\begin{equation}
\delta_{t+1} \leq (1 - \eta\gamma - \frac{1}{n} - \frac{\eta \gamma}{n})\delta_t + \frac{1}{n}\delta_t + \frac{2\eta G}{n},
\end{equation}

\begin{equation}
\delta_{t+1} \leq \delta_t(1 - n\gamma - \frac{1}{n} - \frac{\eta\gamma}{n} + \frac{1}{n}) + \frac{2\eta G}{n},
\end{equation}

\begin{equation}
\delta_{t+1} \leq (1 - \frac{\gamma}{\eta}(\frac{n-1}{n}))\delta_t + \frac{2\eta G}{n}.
\end{equation}

\noindent
This inequality expresses a linear recurrence of the form $\delta_t \leq a\delta_t + b$ with $\delta_0 = 0$, where $a = 1 - \frac{\gamma}{\eta}(\frac{n-1}{n})$ and $b = \frac{2\eta G}{n}$. Expressed as a summation for $\delta_{T}$, this linear recurrence is

\begin{equation}
\delta_{T} \leq a^T\delta_0 + b \sum_{k=0}^{T-1}a^k.
\end{equation}

\noindent
$\sum_{k=0}^{T-1}a^k$ is a geometric series, so $\sum_{k=0}^{T-1}a^k = \frac{1-a^T}{1-a}$. Moreover, since $0 < a < 1$, $a^T \rightarrow 0$ as $T \rightarrow \infty$. This implies

\begin{equation}
\delta_T \leq b(\frac{1}{1-a}),
\end{equation}

\noindent
which yields

\begin{equation}
\delta_T \leq \frac{2 \eta G}{n}(\frac{1}{\eta \gamma}) = \frac{2G}{\gamma n}.
\end{equation}

\noindent
Since the loss function is $G$-Lipschitz, we know $|f(w) - f(w')| \leq G||w - w'||$. Thus, $|f(w)_T - f(w')_T| \leq G\delta_T \leq G(\frac{2G}{\gamma n}) = \frac{2G^2}{\gamma n}$. Thus, $\delta_T \leq \frac{2G^2}{\gamma n}$. Moreover, since $\delta_T \leq \epsilon = O(\frac{1}{n})$, it follows that SGD under $\gamma$-strong convexity and $\beta$-smoothness, with constant $\eta \leq \frac{\gamma}{\beta^2}$ ran over $T$ iterations, is uniformly stable. 
\end{proof}

\end{document}
